\reportsection{03}{Chassis}{The Chassis: Extract, Patch by Regex, Repack}

\unithead{Mechanism}

\reppar{Before any behavior patch runs, the chassis unpacks the upstream archive: \code{asar extract app.asar app.asar.contents}, per-feature patches, then \code{asar pack ... --unpack '**/*.node'} in \code{finalize\_app\_asar} (\code{scripts/staging/electron.sh}). The \code{--unpack} flag is load-bearing: Electron's asar-to-\code{.unpacked} redirect fires only when the manifest entry is annotated as unpacked; without it, \code{require()} of native helpers returns \code{MODULE\_NOT\_FOUND}. Around the extract/repack cycle, \code{patch\_app\_asar} (\code{scripts/patches/app-asar.sh}) performs the \textbf{layout normalization} the Windows-extracted tree demanded: it rewrites \code{package.json} so \code{main} points at \code{frame-fix-entry.js} (a two-line shim that loads \code{frame-fix-wrapper.js} before the original entry point), sets \code{desktopName} to \code{claude-desktop.desktop} (or the \code{io.github.aaddrick} ID for AppImage, since Electron derives the Wayland \code{app\_id} from this field, \#561, \#562), copies the locale JSONs from beside the asar into \code{resources/i18n/} where the app actually resolves them (\#23, en-US.json not found), and copies \code{Tray*} icons into the asar so both packaged and unpackaged (Nix, \code{isPackaged=false}) lookups succeed. It also fails fast if upstream ever renames \code{productName}: Electron ignores \code{--class=} and derives X11 \code{WM\_CLASS} from \code{productName}, so \code{readonly WM\_CLASS='Claude'} in \code{build.sh} is the single source of truth and a mismatch aborts the build.}

\reppar{Because upstream re-minifies every release, the chassis also carries \code{extract\_electron\_variable} in \code{scripts/patches/\_common.sh}, which recovers the electron module variable at build time:}

\begin{codeblock}
electron_var=$(grep -oP '[$\w]+(?=\s*=\s*require\("electron"\))' \
    "$index_js" | head -1)
\end{codeblock}

\reppar{with a fallback anchor on \code{new X.Tray} and a hard exit if both fail. Its sibling \code{fix\_native\_theme\_references} auto-repairs a recurring upstream minifier bug where the bundle references \code{nativeTheme} through the wrong alias.}

\unithead{Origin}

\reppar{The chassis predates every feature patch: extract, tray-icon copy, and repack are already present in \code{d8f4bbe}, the repository's December 2024 initial commit (\code{build-deb.sh}). Everything layered on since is a normalization or repair forced by the Windows-extracted tree or by upstream re-minification: the i18n copy from Coldsewoo's \#23 report, the \code{main} rewrite from speleoalex's \#127 (windows appeared undecorated on Linux window managers), dynamic identifier extraction from \#220 (fixing \#218/\#219), and the productName/WM\_CLASS fail-fast from \#655.}

\begin{chronology}
2024-12 & \code{d8f4bbe} & Origin: initial \code{build-deb.sh} commit already contains the asar extract/repack scaffolding and an in-asar tray-icon copy. \\
2025-03 & \#23, \#26 \code{466705d}; \#25 \code{ee7e4bc} & i18n locale JSON copy into \code{resources/i18n/} added, motivated by Coldsewoo's \#23 (\code{en-US.json} not found inside the asar); \#26 landed the surviving \code{*-*.json} glob first, \#25 merged later as the documented fix (the \#23 link for \#26 is inferred from identical code, not recorded). \\
2025-11 & \#127, \code{7882635} & package.json \code{main} redirected through \code{frame-fix-entry.js} (speleoalex): Claude Desktop windows appeared without borders or decorations on Linux window managers. \\
2025-11 & \#122, \code{47c1889} & Tray icons pulled back out of the asar for a filesystem-only copy: "Tray icons must be in filesystem (not inside asar) for Electron Tray API to access them." \\
2026-02 & \#218/\#219, \code{7bd2f38} (PR \#220) & Anchor rot 1: upstream v1.1.2321 re-minified the electron variable \code{oe} to \code{Ae}, breaking every hardcoded tray sed; hardcoded lookups replaced with \code{extract\_electron\_variable} plus \code{fix\_native\_theme\_references}. \\
2026-02 & \#252, \code{546f845} (PR \#253) & Anchor rot 2: v1.1.4088 minified the electron variable to \code{\$e}; \code{\textbackslash w} does not match \code{\$}, so extraction captured \code{e} and seds spliced code mid-name, a launch-time SyntaxError; \code{electron\_var\_re} introduced for regex-escaped sed use. \\
2026-02 & \code{573f052} (PR \#266) & Tray icons copied back into the asar so unpackaged builds (Nix, \code{isPackaged=false}) can find them, reversing \#122's filesystem-only move. \\
2026-04 & \#421, \code{3150477} & \code{--unpack '**/*.node'} added to the repack after node-pty \code{require()} returned \code{MODULE\_NOT\_FOUND}: the asar-to-\code{.unpacked} redirect needs the manifest entry annotated as unpacked. \\
2026-04 & \code{ff4821e} & The 2124-line \code{build.sh} split into modules, creating \code{scripts/patches/\_common.sh}, \code{scripts/patches/app-asar.sh}, and \code{scripts/staging/\{electron,locales\}.sh} (pure move, function bodies verbatim). \\
2026-05 & \#561/\#562, \code{5a98854} & \code{desktopName} edit added (AppImage-specific ID): Electron derives the Wayland \code{app\_id} from this field, and KDE Wayland grouped the pinned launcher separately from the running window. \\
2026-05 & \#644, \code{b40441c} & Anchor rot 3: \code{\textbackslash w} silently truncated mid-\$ names like \code{i\$A}; \code{[\$\textbackslash w]+} hardened repo-wide (3 sites in \code{\_common.sh}), the same \$-capture trap having recurred in \#421 and \#555 before this convention stuck. \\
2026-05 & \#648, \code{76a5a21} & WM\_CLASS and StartupWMClass aligned to \code{claude-desktop} across all formats: the wrong direction, superseded two days later. \\
2026-05 & \#652/\#655, \code{e7e6475}/\code{73c9b8f} & Reversal and centralization: Electron ignores \code{--class=} and derives \code{WM\_CLASS} from \code{productName} (\code{Claude}); users confirmed via \code{/proc} cmdline plus \code{xprop} that the flag was silently ignored. Added \code{readonly WM\_CLASS='Claude'} and the productName fail-fast. \\
\end{chronology}

\methodbanner{\textbf{Fate under the rebase.} \brework{} The chassis is reworked, not deleted. With the official \code{.deb} a conventional electron-forge Linux tree, most normalization has nothing left to fix: \code{\_common.sh} is gone (\code{grep -rn electron\_var scripts/} returns nothing; both survivors self-anchor on literals), the frame-fix \code{main} edit is deleted ("the only \code{frame:!1} sites are the Quick Entry popup and two transparent overlay windows"), the \code{desktopName} edit and the tray-icon copy are deleted (upstream ships \code{TrayIconLinux(-Dark).png} natively), and the i18n copy is safely dead: byte-checked against the pinned official 1.17377.2 deb, locale JSONs ship loose at \code{usr/lib/claude-desktop/resources/*.json} (plus \code{resources/ion-dist/i18n/}). What remains on \code{rebase/official-deb} is a thin orchestrator with \code{active\_patches=(patch\_quick\_window patch\_org\_plugins\_path)}; an empty array means the official \code{app.asar} ships \textbf{byte-identical}, no extract, no repack. When patches are active, the \code{--unpack} expression is derived from the shipped \code{app.asar.unpacked} tree, folded into a single brace glob, and guarded by a set-equality check that hard-fails if the repacked unpacked set diverges. The one piece that survives upgraded is the productName/WM\_CLASS fail-fast, now run unconditionally via \code{asar extract-file} so the assertion holds even on a patch-zero build that never unpacks the asar.\\[3pt]\textbf{Affects.} All Linux; packaging-layer normalization, format-general (deb/AppImage/Nix). The Wayland \code{app\_id} derivation (\#561/\#562) is the one environment-specific edge.}
